% ============================================================
%  Appendix C — Deep-Learning Foundations
%
%  DRAFT (AI-assisted, June 2026). This appendix grounds the
%  deep-learning machinery applied in Chapter 2 in the theory and
%  reading list of the MSc Deep Learning course. It is written as
%  scaffolding: revise into your own voice, trim anything the examiner
%  would consider textbook padding, and keep only what earns its place.
%  Every claim is cited; cross-references point back to the methods in
%  Chapter 2 so the appendix documents *why* each choice is principled
%  rather than restating it.
% ============================================================

\chapter{Deep-Learning Foundations}
\label{app:dlfoundations}
% SUGGESTED TRIM: KEEP but tighten. Opening framing is fine; trim sentences that only announce "this appendix is background" rather than doing work.

The methodology of Chapter~\ref{chap:methodology} applies a sequence of
deep-learning tools --- a multilayer network trained by stochastic
gradient descent, adaptive optimisation, dropout, $\ell_2$
regularisation, early stopping, and ensembling --- as established
technique, without pausing to justify them. This appendix supplies
that justification, drawing on the standard references
\citep{goodfellow2016,hastie2009}. It is deliberately background; the
intent is to record the theoretical results that make each modelling
choice in Chapter~\ref{chap:methodology} principled rather than
heuristic, and to place the study in the lineage of machine learning
applied to credit and asset-pricing problems.

\section{Function classes and universal approximation}
\label{sec:dlf-approx}
% SUGGESTED TRIM: KEEP (defends a Ch.2 choice). Justifies the deep hypothesis class over the logit; compress the textbook universal-approximation exposition, retain the "logit = zero-hidden-layer boundary case" point.

The estimand \eqref{eq:object} is a conditional distribution over next
month's state, and the modelling question is how flexibly a candidate
$p_\theta(\cdot\mid x)$ is allowed to depend on the covariates $x$. The
multinomial logit \eqref{eq:logit} confines that dependence to a
log-linear form; the deep network \eqref{eq:mlp} does not. The
theoretical warrant for the latter is the universal approximation
theorem: a feedforward network with a single hidden layer of
sufficiently many units and a non-polynomial activation can approximate
any continuous function on a compact set to arbitrary accuracy
\citep{cybenko1989,hornik1991}. The result certifies that the
hypothesis class of \eqref{eq:mlp} contains, in the limit of width, the
true transition surface whose nonlinearity \citet{sirignano2021} claim
is economically material --- so a deep model's empirical edge over the
logit reflects a representable feature of the data, not an artefact.
The theorem is, however, existential and says nothing about whether a
finite network can be \emph{found} by training, nor how much data the
search needs; depth (searched over $L\in\{1,3,5,7\}$ in
Section~\ref{subsec:meth-nn}) is the practically more economical route
to the same expressive power, building the interaction hierarchy ---
incentive$\times$credit-score and the like --- through composition
rather than width. The logit is then exactly the zero-hidden-layer
boundary case of the same class (Section~\ref{subsec:meth-logit}),
which is what makes the model hierarchy a controlled comparison.

\section{Training: backpropagation and stochastic gradient descent}
\label{sec:dlf-sgd}
% SUGGESTED TRIM: KEEP core (defends the two-level shuffle + Horvitz-Thompson weights = unbiased minibatch gradient). Compress the backprop / Robbins-Monro / Bottou history.

Fitting \eqref{eq:mlp} means minimising the weighted negative
log-likelihood \eqref{eq:weightednll} over the parameters $\theta$.
The gradient $\nabla_\theta\widehat{\mathcal L}$ is computed by
backpropagation --- reverse-mode automatic differentiation, the
reverse accumulation of the chain rule through the network's
computational graph \citep{rumelhart1986} --- which yields all partial
derivatives at the cost of a constant multiple of one forward
evaluation, and is what PyTorch's autograd implements. Because the
objective is an average over $\sim$$10^{9}$ loan-months, full-batch
gradients are infeasible, and the parameters are updated on minibatch
estimates: stochastic gradient descent, whose convergence under
decreasing step sizes traces to the stochastic-approximation theory of
\citet{robbins1951} and is treated for the modern large-scale setting
by \citet{bottou2018}. SGD's guarantees rest on minibatch gradients
being \emph{unbiased} estimates of the full-data gradient. Two design
choices in Chapter~\ref{chap:methodology} exist precisely to protect
that property: the loan-keyed two-level shuffle of
Section~\ref{sec:meth-shards}, which keeps each minibatch an
approximately uniform cross-section despite vintage-grouped storage,
and the Horvitz--Thompson inverse-probability weights in
\eqref{eq:weightednll}, which keep the thinned-sample loss --- and
hence its gradient --- an unbiased estimate of the full-sample loss.

\section{Adaptive optimisers}
\label{sec:dlf-adam}
% SUGGESTED TRIM: KEEP (defends a Ch.2 choice: Adam used identically for the logit control, so the comparison varies architecture alone). Already short.

Plain SGD with a single global step size is sensitive to feature
scaling and curvature anisotropy, both severe here given embeddings and
heavy-tailed balances. The training runs therefore use Adam
\citep{kingma2015}, which maintains per-parameter exponential moving
averages of the gradient and its square and rescales each update by the
latter, combining the per-coordinate adaptivity of RMSprop
\citep{tieleman2012} with momentum. This is the optimiser referenced in
the loan-level model plan (Adam, learning rate $\sim$$10^{-3}$ with
decay-on-plateau), and it is used identically for the logit control of
Section~\ref{subsec:meth-logit}, so that the network-versus-logit
comparison varies architecture alone and not the optimisation
procedure.

\section{Regularisation, ensembling, and normalisation}
\label{sec:dlf-reg}
% SUGGESTED TRIM: MIXED. KEEP the batch/layer-norm OMISSION defence (justifies a Ch.2 design choice); COMPRESS the textbook ridge / dropout / bagging recap.

A width-200, several-layer network has ample capacity to overfit rare
credit events, and Section~\ref{subsec:meth-nn} deploys four classical
controls. The $\ell_2$ penalty (weight decay) is ridge regularisation
in the sense of \citet{hastie2009}, shrinking parameters toward zero and
trading a little bias for reduced variance. Dropout
\citep{dropout} randomly zeroes hidden units during training, which
acts as an implicit ensemble over subnetworks and a stochastic
regulariser. Early stopping on the window's validation likelihood halts
the fit before it specialises to training noise --- a form of
regularisation through optimisation budget. The headline ensemble, an
average of $K{=}8$ independently initialised networks, is bagging-style
variance reduction \citep{hastie2009}: averaging quasi-independent
predictors lowers the variance component of generalisation error, which
for high-capacity models fit by a stochastic procedure is substantial.

Normalisation layers --- batch normalisation \citep{ioffe2015} and layer
normalisation \citep{ba2016} --- are part of the standard deep-learning
toolkit and are covered in the course, but the loan-level model does
\emph{not} use them. For this tabular network the per-window
standardisation of continuous covariates described in
Section~\ref{sec:meth-features} supplies the input conditioning that
batch normalisation would otherwise provide, and avoiding
batch-statistic layers keeps the strictly out-of-sample design of
Section~\ref{sec:meth-design} clean (no train-time batch statistics
leaking across the temporal split). This is a deliberate omission, not
an oversight, and is noted here for completeness.

\section{Architectures beyond the feedforward network}
\label{sec:dlf-arch}
% SUGGESTED TRIM: KEEP the "why a plain feedforward net" argument (defends the Ch.2 Markov framing); COMPRESS the ResNet/LSTM/Transformer/GAN survey to a sentence.

The course surveys architectures this study does not use ---
residual networks \citep{he2016}, recurrent networks such as the LSTM
\citep{hochreiter1997}, attention-based transformers
\citep{vaswani2017}, and generative adversarial networks
\citep{goodfellow2014gan} --- and it is worth stating why a plain
feedforward network is the right tool here. The modelling object
\eqref{eq:object} is a \emph{one-month} conditional transition given
the current state and covariates; the Markov framing of
Section~\ref{sec:meth-model} already encodes the loan's dynamics
through the state feature and loan age, so the temporal dependence that
would motivate a recurrent or attention model is handled by the
state-space construction rather than by the network. Residual
connections and normalisation address optimisation pathologies of very
deep networks that do not arise at depths $L\le 7$. The feedforward
choice is thus a consequence of the problem's structure, made against
the full menu of alternatives rather than by default.

\section{Machine learning in credit and asset pricing}
\label{sec:dlf-lineage}
% SUGGESTED TRIM: KEEP (literature positioning, not padding). The Gu-Kelly-Xiu "levels not ranks" link to the pool-level pricing exercise earns its place.

The study sits in an established line of work applying flexible
machine learning to financial prediction. \citet{sirignano2021} ---
the paper this dissertation tests on a cleaner dataset --- is the
direct antecedent for deep learning of mortgage transition risk.
Earlier, \citet{khandani2010} showed machine-learning models improve
consumer credit-risk forecasts, and \citet{fuster2022} examine the
distributional consequences of machine learning in credit markets,
both relevant to the credit-event side of the transition model.
\citet{gukellyxiu2020} establish the value of nonlinear machine
learning for the levels of expected returns in empirical asset
pricing --- the methodological cousin of the pool-level pricing
exercise in Section~\ref{sec:meth-pool}, where prediction levels (not
merely rankings) are what the valuation consumes. Reinforcement
learning, surveyed for finance by \citet{bai2025} and developed in
\citet{suttonbarto2018}, lies outside this study's supervised framing
but is the natural setting for the sequential decision problems
(prepayment-option exercise, dynamic hedging of a pool position) that
the transition model's outputs would feed.

\section{Course-to-dissertation crosswalk}
\label{sec:dlf-crosswalk}
% SUGGESTED TRIM: KEEP (crosswalk table, explicitly retained). The course-to-dissertation map is the appendix's load-bearing exhibit.

Table~\ref{tab:dlcrosswalk} maps each topic from the Deep Learning
course to where it is used in this dissertation and to the supporting
reading.

\begin{table}[htbp]
\centering
\small
\renewcommand{\arraystretch}{1.2}
\begin{tabularx}{\textwidth}{@{}lXl@{}}
\toprule
\textbf{Course topic} & \textbf{Where it is used in this dissertation} & \textbf{Reading} \\
\midrule
Fully-connected networks &
  Loan-level transition model, Eq.~\eqref{eq:mlp}, Sec.~\ref{subsec:meth-nn} &
  \citealp{goodfellow2016} \\
Universal approximation &
  Justifies the deep hypothesis class over the logit, App.~\ref{sec:dlf-approx} &
  \citealp{cybenko1989,hornik1991} \\
Backpropagation / autodiff &
  Gradient of the weighted NLL, Eq.~\eqref{eq:weightednll} &
  \citealp{rumelhart1986} \\
SGD / minibatch SGD &
  Two-level shuffle over shards, Sec.~\ref{sec:meth-shards} &
  \citealp{robbins1951,bottou2018} \\
Adam / RMSprop &
  Optimiser for all fits, Sec.~\ref{subsec:meth-nn} &
  \citealp{kingma2015,tieleman2012} \\
$\ell_2$, dropout, ensembles &
  Three regularisers + $K{=}8$ ensemble, Sec.~\ref{subsec:meth-nn} &
  \citealp{hastie2009,dropout} \\
Batch / layer normalisation &
  Surveyed; not used (input standardisation instead), App.~\ref{sec:dlf-reg} &
  \citealp{ioffe2015,ba2016} \\
Residual / recurrent / attention &
  Surveyed alternatives not required by the Markov framing, App.~\ref{sec:dlf-arch} &
  \citealp{he2016,hochreiter1997,vaswani2017} \\
ML in finance &
  Lineage of the study, App.~\ref{sec:dlf-lineage} &
  \citealp{sirignano2021,gukellyxiu2020,khandani2010,fuster2022} \\
\bottomrule
\end{tabularx}
\caption[Deep Learning course topics mapped to dissertation methods]{Each
Deep Learning course topic, the place it is used in this dissertation,
and the supporting reading.}
\label{tab:dlcrosswalk}
\end{table}
